%!TEX root = ../main.tex
% APÉNDICE B — Entrevistas a especialistas: guiones
% Los guiones completos se referencian desde el cap. 3 (3.4).

\chapter{Entrevistas a especialistas}\label{apx:entrevista}

\begin{guia}
Apéndice por entrevista: contexto (fecha, medio, duración, objetivo),
datos de la persona entrevistada (formación, cargo, trayectoria, con su
consentimiento) y la transcripción organizada por temas. Las
conclusiones van en el cuerpo (capítulos \ref{cap:marco-teorico} y
\ref{cap:resultados}); acá va la evidencia.
\end{guia}

Como se detalla en el capítulo~\ref{cap:marco-teorico}
(sección~\ref{sec:entrevistas}), se definieron dos perfiles de
especialista complementarios: un docente o gestor/a institucional y un/a
profesional del campo psicopedagógico o del vínculo y la seguridad con
menores. En cada caso se documenta el contexto, la ficha de la persona
entrevistada y la transcripción por ejes temáticos.

\textbf{Estado: [PENDIENTE]} -- el cronograma del
capítulo~\ref{cap:metodologia} prevé su realización durante el Sprint 3.
Hasta entonces los guiones aquí transcritos constituyen la evidencia del
diseño metodológico; los \completar{completar} marcan los datos que se
completarán con la realización efectiva de cada entrevista.

\section{Perfil 1: Docente o gestor/a institucional}

\subsection{Contexto de la entrevista}

\completar{Fecha, modalidad (presencial/remota), medio y duración.}

\subsection{Datos de la persona entrevistada}

\completar{Formación, cargo actual, experiencia relevante.}

\subsection{Guion de entrevista}

\emph{Eje 1 -- apoyo escolar y repitencia.}

\begin{enumerate}
    \item ¿Qué rol juega el apoyo escolar personalizado en la prevención
          de la repitencia en los niveles primario y secundario?
    \item ¿Qué diferencias observa entre el acceso a tutorías en zonas
          urbanas y rurales del país?
\end{enumerate}

\emph{Eje 2 -- tutoría virtual pedagógica y de seguridad.}

\begin{enumerate}
    \item Desde lo pedagógico, ¿qué recaudos considera necesarios para que
          una tutoría virtual 1 a 1 sea efectiva con un menor?
    \item Desde la seguridad, ¿qué condiciones mínimas debería garantizar
          una plataforma que conecta adultos desconocidos con menores en
          un entorno remoto?
    \item ¿Cómo observaría la articulación de esta plataforma con las
          instituciones educativas (escuelas, docentes, equipos de
          orientación)?
\end{enumerate}

\subsection{Transcripción de la entrevista}

\completar{Transcripción organizada por ejes temáticos.}

\section{Perfil 2: Psicopedagogía / vínculo y seguridad con menores}

\subsection{Contexto de la entrevista}

\completar{Fecha, modalidad (presencial/remota), medio y duración.}

\subsection{Datos de la persona entrevistada}

\completar{Formación, cargo actual, experiencia relevante.}

\subsection{Guion de entrevista}

\emph{Eje 1 -- vínculo y acompañamiento.}

\begin{enumerate}
    \item ¿Qué impacto vincular observa en un acompañamiento personalizado
          1 a 1 para niños, niñas y adolescentes?
    \item ¿Qué señales de alerta de un menor conviene monitorear en un
          entorno remoto de interacción con un adulto no conocido?
    \item ¿Qué rol recomienda para el adulto responsable dentro de la
          plataforma?
\end{enumerate}

\emph{Eje 2 -- reputación y diseño seguro.}

\begin{enumerate}
    \item ¿Qué errores evitaría en el diseño de un sistema de reputación
          bidireccional que involucra a menores?
    \item ¿Qué balance propone entre la protección del menor y el respeto a
          su privacidad al revisar sesiones o interacciones?
    \item ¿Qué recomendaciones daría para el manejo de denuncias y
          situaciones de conflicto entre tutores y familias?
\end{enumerate}

\subsection{Transcripción de la entrevista}

\completar{Transcripción organizada por ejes temáticos.}